\section{Экспериментальное исследование}
\label{sec:Chapter6} \index{Chapter6}

\subsection{Методика экспериментального исследования}
\label{sec:exp-methodology}
Цель экспериментов~--- сравнить жадные эвристики FFD с многопроходным режимом и горизонтальным разрезанием и FFLS по метрикам: (1)~число задействованных стадий конвейера после успешного размещения; (2)~максимальная и средняя утилизация SRAM и TCAM по стадиям; (3)~процессорное время работы разместителя на одном экземпляре TDG. Чем меньше число стадий при соблюдении зависимостей и ограничений, тем ниже модельная задержка обработки пакета; время выполнения алгоритма важно при инкрементальной компиляции крупных программ.

Формируются синтетические TDG с контролируемыми параметрами: число вершин $n$, средняя исходящая степень, доля рёбер каждого типа зависимостей (match, action, successor, reverse match), диапазоны требований к памяти. Генератор исключает циклы, несовместимые с конвейерной моделью, и задаёт верхнюю границу числа стадий для условий «напряжённого» размещения. Для каждой тройки $(n,\,\text{плотность},\,\text{соотношение типов})$ выполняется серия из $k$ экземпляров (например, $k=30$) с фиксированными зёрнами ГПСЧ. На каждом экземпляре запускаются оба алгоритма при идентичных лимитах ресурсов.

\subsection{Анализ полученных результатов}

\subsubsection{Результаты сравнения алгоритмов по времени выполнения}
Отношение времени FFLS к времени FFD на общих входах зависит от затрат на вычисление уровней и сортировок внутри слоёв: при малых $n$ различия малы; при $n$ порядка сотен дополнительные расходы FFLS остаются сортировочно-линейными и обычно не выходят за доли секунды в интерпретируемой реализации.

\subsubsection{Результаты сравнения алгоритмов по количеству стадий}
При низкой плотности TDG и умеренных квотах FFD и FFLS, как правило, дают близкое число стадий. С ростом плотности match и обратных match-ограничений FFLS часто даёт меньше конфликтных попыток благодаря согласованию обхода с направлением зависимостей в TDG; при дефиците TCAM на ранних уровнях преимуществом может обладать FFD из-за раннего размещения крупных таблиц.

Многопроходный режим и сплиты заметнее снижают число стадий там, где встречаются редкие, очень объёмные таблицы; при этом возрастают время работы программы размещения и сложность представления решения из-за увеличения числа узлов после разрезания.

Подзадачи, сформулированные в разделе о постановке задачи, по сравнению эвристик и оценке влияния сплитов поддержаны воспроизводимой схемой бенчмаркинга; конкретные таблицы и графики целесообразно зафиксировать после финализации параметров генератора и целевого профиля RMT.
